% Pursuing Stacks (cotes 134-2 to 134-8): archive cote 134-6, site page 23 (PDF page 24), typescript page 342.
% typescript. Draft compared: ps5.tex:2510-2569.
% AI-transcribed by Claude Opus 5.5 (claude-opus-5-5, adaptive thinking, effort high), Message Batches API, batch ps3 (msgbatch_01XdFmuVVBAcgc7EFxdtCqsj), 2026-11-12; not yet checked by eye.
\scanpage{134-6/23}
\mspage{342}

ons about seven years ago.  Replacing $\underline{Q}$ by a more general ring,  this
suggests that \uline{there} \xout{shoul} \uline{might exist a notion of\hand{``}homotopy types\hand{''}over}
\uline{any ground ring}  k  -  and a corresponding notion of ground ring exten-%
sion for homotopy types.  For abelianizations of homotopy types,  this
is particularly "obvious" ,  as  being just the functorial dependance
of the derived category  D.(k)  with respect to the ground ring k,
corresponding to ring extension in a chain complex.  For a \xout{few} week or
two I played around with this idea,  which on the semisimplicial level
tied in with expressing homotopy types of some simple spaces (\over{\unsure{as}}{such as}
standard K($\pi$,n) spaces and fibrations between these) in terms of some
\inserted[typed]{simple} \uline{semisimplicial} \uline{schemes} \inserted[typed]{(affine and} of finite type over Spec ($\underline{Z}$ )),  by taking $\underline{Z}$-%
valued points of these\over{,}{\unsure{;}} \deleted{whereas} ring extension  $\underline{Z} \longrightarrow$ k  was interpre-%
ted in the scheme-theoretic \uline{sense.}\margin{$\rfloor\!\lceil$}I did'nt go on very long,  as\inserted{soon after}I was
\deleted{then} taken by personal matters and never took up the matter later  -
and maybe it was an altogether unrealistic or silly attempt.  If I re-%
member it right,  the idea lurking was \xout{that there} something of this
kind,  that \uline{there was a functor from} (Hot)\inserted{to}(if not an equivalence of
(Hot) with ...)  \xout{to} \uline{a suitable derived category of} \deleted{the}\inserted{some} \uline{category of}
\uline{semisimplicial schemes  over} Spec($\underline{Z}$) ,  and that the base change intui-%
tion,  as suggested by the abelianized theory or by the subtler "divided
\inserted{power} De Rham theory",  would reflect in naive base change  Z $\longrightarrow$ k  for
\uline{schemes.}\margin{$\rfloor\!\lceil$}I was then looking mainly at 1-connected structures,  but there
was an idea too that nilpotent fundamental groups might fit into the
picture,  with the hope that such a group (under suitable restrictions\hand{,}
\deleted{such as} finite presentation and torsion freeness say)  could be \xout{expresse}
expressed in a canonical way in terms of an affine\inserted{nilpotent}group scheme \xout{over}
of finite type over Spec($\underline{Z}$),  by taking the integral points of the
latter.  It seems (if I remember right)  that this is not quite true\xout{,}
though  -  that one could'nt hope for much better than getting a nilpo-%
tent algebraic group scheme \uline{over $\underline{Q}$}  -  and that one would recover the
discrete group one started with only\hand{\textbackslash}"up to commensurability" . \inserted{Possibly}
\deleted{thus getting possibly}\inserted{there may be} \uline{an equivalence between} \deleted{the} \uline{localization of the}
\uline{category of nilpotent groups as above} (\uline{with respect to monomorphisms}
\uline{with image of finite index})  \uline{and}\inserted[typed]{affine} \uline{nilpotent connected algebraic group}
\uline{schemes over $\underline{Q}$} ,  or equivalently,  group schemes whose underlying
scheme is isomorphic to standard affine space.

\margin{23}
